
\begin{definition}\textbf{Integral de trayectoria de campos escalares.}
    Sea $\boldsymbol{\sigma}:I=[a,b]\to\Rn{n}$ una trayectoria de clase $\mathcal{C}^1(\interior{I})$ y sea $f$ un campo escalar tal que este bien definida y sea continua   la funci\'on $f\circ\boldsymbol{\sigma}$.  Se define la \emph{integral de trayectoria}, o \emph{integral de} $f$ \emph{a lo largo de la trayectoria} $\boldsymbol{\sigma}$, y se denota por $\int_{\boldsymbol{\sigma}}f\:ds$, a
    \[
        \int_{\boldsymbol{\sigma}}f \, ds=\int_a^b (f\circ\boldsymbol{\sigma})(t)  \|\boldsymbol{\sigma}'(t)\| \, dt.
    \]
\end{definition}

 \begin{example} 
 Sea $\boldsymbol{\sigma} : [0, 3\pi] \to \mathbb{R}^2$ / 
$\boldsymbol{\sigma}(t) = (\cos t,\, \sen t)$. Calcular $\displaystyle\int_{\boldsymbol{\sigma}} f\, ds$,
donde $f(x,y) = xy$.
\end{example}

\noindent\textbf{Solución:} Calculamos primero el campo evaluado en la trayectoria y la norma del vector tangente:
\begin{align*}
f(\boldsymbol{\sigma}(t)) &= \cos t \cdot \sin t, \\[0.2cm]
\boldsymbol{\sigma}'(t) &= (-\sin t,\, \cos t) \implies \|\boldsymbol{\sigma}'(t)\| = \sqrt{\sin^2 t + \cos^2 t} = 1.
\end{align*}
Planteamos la integral sustituyendo los elementos obtenidos:
\[
    \int_{\boldsymbol{\sigma}} f\, ds
    = \int_0^{3\pi} f(\boldsymbol{\sigma}(t))\,\|\boldsymbol{\sigma}'(t)\|\, dt
    = \int_0^{3\pi} \cos t \cdot \sen t \, dt.
\]
Resolvemos:

\begin{align*}
\int_0^{3\pi} \cos t \cdot \sin t\, dt &= \int_0^{3\pi} \frac{\sin(2t)}{2}\, dt \\[0.2cm]
&= \left[ -\frac{\cos(2t)}{4} \right]_0^{3\pi} \\[0.2cm]
&= -\frac{\cos(6\pi)}{4} + \frac{\cos(0)}{4} \\[0.2cm]
&= -\frac{1}{4} + \frac{1}{4} = 0.
\end{align*}
Por lo tanto, se concluye que $\int_{\boldsymbol{\sigma}} f\, ds = 0$. \hfill $\blacksquare$ 


% Si $\boldsymbol{\sigma}(t)$ es de clase $\mathcal{C}^1$ a trozos, se define $\int_{\boldsymbol{\sigma}}f\:ds$  de la siguiente manera: 

 
\begin{definition}\label{def:int_linea_esca_a_trozos}  Sea $\boldsymbol{\sigma}:I=[a,b]\to\Rn{n}$ una trayectoria de clase $\mathcal{C}^1$ a trozos   y  sea $f$ un campo escalar tal que este bien definida y sea continua   la funci\'on composici\'on $f\circ\boldsymbol{\sigma}$.  Se define la \emph{integral de trayectoria}  o \emph{integral de} $f$ \emph{a lo largo de la trayectoria} $\boldsymbol{\sigma}$, y se  nota $\int_{\boldsymbol{\sigma}}f\:ds$, a
  
    \[
       \int_{\boldsymbol{\sigma}}f \, ds =\sum_{i=0}^{m-1}\int_{I_i}  (f\circ\boldsymbol{\sigma})(t)\|\boldsymbol{\sigma}'(t)\| \, dt,  
    \]
    donde $ I_i= (t_i, t_{i+1})$,  $\boldsymbol{\sigma}\lvert_{I_i}$ es de clase $\mathcal{C}^1$  y  $\{t_i\}_{i=0}^{m}$  es una partici\'on de $I$.
\end{definition}

\begin{example} 

Sea $\boldsymbol{\sigma} : [0, 2] \to \mathbb{R}^2$ la trayectoria de clase $\mathcal{C}^1$ a trozos definida por 
\[ 
\boldsymbol{\sigma}(t) = \begin{cases} (t,\, 0) & \text{si } t \in, \\ (1,\, t - 1) & \text{si } t \in. \end{cases} 
\] 
Para  $f(x, y) = x + y$,   calcular $\displaystyle\int_{\boldsymbol{\sigma}} f\, ds$.
\end{example}


\noindent\textbf{Solución:} Dado que la trayectoria está definida a trozos, dividimos el cálculo en dos subintervalos utilizando la partición $\{0, 1, 2\}$ del intervalo $[0,2]$,  analizando cada tramo de forma independiente, seg\'un la  definición ~\ref{def:int_linea_esca_a_trozos}.

\noindent\textbf{Tramo 1 ($t \in [0, 1]$):} Denotamos este primer tramo como $\boldsymbol{\sigma}_1(t) = (t, 0)$ con $t \in [0,1]$. 
\begin{enumerate}
    \item Su vector tangente es $\boldsymbol{\sigma}_1'(t) = (1, 0)$,  por lo que su norma vale $\|\boldsymbol{\sigma}_1'(t)\| = 1$.
    \item Evaluamos el campo escalar sobre este tramo, obteniendo $ (f \circ \boldsymbol{\sigma}_1) (t)  = t + 0 = t$.
\end{enumerate}

\noindent\textbf{Tramo 2 ($t \in [1, 2]$):} Denotamos el segundo tramo como $\boldsymbol{\sigma}_2(t) = (1, t - 1)$  con $t \in [1,2]$. 
\begin{enumerate}
    \item Su vector tangente es $\boldsymbol{\sigma}_2'(t) = (0, 1)$, por lo que su norma vale $\|\boldsymbol{\sigma}_2'(t)\| = 1$.
    \item Evaluamos el campo escalar sobre este tramo, obteniendo $ (f \circ \boldsymbol{\sigma}_2) (t) = 1 + (t - 1) = t$.
\end{enumerate}


\noindent Finalmente, sumamos las integrales de ambos tramos:
\begin{align*}
\int_{\boldsymbol{\sigma}} f\, ds &= \int_{0}^{1}   (f \circ \boldsymbol{\sigma}_1) (t)  \|\boldsymbol{\sigma}_1'(t)\| \, dt + \int_{1}^{2}  (f \circ \boldsymbol{\sigma}_2) (t)  \|\boldsymbol{\sigma}_2'(t)\| \, dt \\[0.2cm]
&= \int_{0}^{1} t \cdot 1 \, dt + \int_{1}^{2} t \cdot 1 \, dt  = \left[ \frac{t^2}{2} \right]_{0}^{1} + \left[ \frac{t^2}{2} \right]_{1}^{2} \\[0.2cm]
&= \left( \frac{1}{2} - 0 \right) + \left( \frac{4}{2} - \frac{1}{2} \right) = 2.
\end{align*}
Por lo tanto, obtenemos que $\int_{\boldsymbol{\sigma}} f\, ds =  2$  \hfill $\blacksquare$






